\begin{claim}{C.65}{Proved}{Principal-kernel boundary comparison and controlled planar fillings}
Let $f_t,G$ be the original plane loop and its generated
group, and let $O$ be an open $G$-orbit with one-dimensional
full local centralizer. Retain the actual principal quotient
of C.33,
\[
 \pi_0:O\longrightarrow B_0=O/C_O^0,
 \qquad
 \beta:G\longrightarrow\operatorname{Homeo}(B_0).
                                                               \tag{C.65.1}
\]
Let $N$ consist of endpoints $h_1$ of paths
\[
\begin{gathered}
 h_s=\prod_{j=1}^{m}f_{a_j(s)}^{\epsilon_j},
 \qquad a_j:[0,1]\to[0,1]\ \text{continuous},\\
 \epsilon_j\in\{1,-1\},
 \qquad h_0=\id,\qquad\beta(h_s)=\id\ \text{for every }s.
\end{gathered}
                                                               \tag{C.65.2}
\]
The finite length and signs may depend on the path.
For every nontrivial nonopen $G$-orbit $B\subset\partial O$
whose actual intrinsic restriction image $\rho_B(G)$
is nonmetabelian,
\[
              \rho_B(N)=1,\qquad
 \rho_B:G\longrightarrow\operatorname{Homeo}(B).
                                                               \tag{C.65.3}
\]
In particular two actual finite-word paths $p_s,q_s$
starting at the identity and satisfying
$\beta(p_s)=\beta(q_s)$ for every $s$ have exactly the
same restriction to each such $B$ at every time.

For the full principal kernel $K_0=\ker\beta$, the
image $D_B=\rho_B(K_0)$ is central in $\rho_B(G)$
and is a closed discrete cyclic subgroup of the
intrinsic homeomorphism group of $B$; it is finite
when $B$ is a circle. Give
\[
 E_B=\{(\beta(g),\rho_B(g)):g\in G\}
       \subset\operatorname{Homeo}(B_0)
                    \times\operatorname{Homeo}(B)
\]
the actual induced product topology, and put
$\Delta_B=\{1\}\times D_B$. There are exact maps
\[
 E_B\longrightarrow E_B/\Delta_B
             \longrightarrow\beta(G).
\]
The first map is a covering with deck group
$\Delta_B$; the second is a continuous bijective
homomorphism onto the actual induced image
$\beta(G)$. No continuity of its inverse is asserted.

There is also the following controlled filling fact for
the same planar setting. For any open planar set $V$,
any open Euclidean disk $D$, and any connected component
$U$ of $V\cap D$, the actual inclusion induces an
injection
\[
                        \pi_1(U,x)\longrightarrow\pi_1(V,x)
                    \qquad(x\in U).                    \tag{C.65.4}
\]
Thus a based loop in $U$ null-homotopic in $V$ has a
based disk filling entirely in $U$. Applied to shrinking
null evaluation loops near a fixed frontier point,
this gives fillings inside the original orbit with
vanishing ambient diameter.
\end{claim}

\begin{proof}
\emph{The actual kernel paths and their intrinsic topology.}
Every $g\in G$ commutes with $C_O^0$, by the definition
of the full local centralizer. It therefore induces
the actual base homeomorphism in (C.65.1). This
defines a homomorphism, with no change of quotient.
The set $N$ in (C.65.2) is a subgroup: products
and inverses of two displayed paths are again
finite-word paths in the exact kernel, and
identity letters can be inserted if needed.
Conjugation by a fixed original word preserves
that kernel and adds only finitely many constant
letters. Hence $N$ is normal in $G$.
Every intermediate $h_s$ also belongs to $N$,
by reparametrizing the path on the interval from
$0$ to $s$.

We verify the topology used on $B$. On a circular
orbit the intrinsic and embedded circle topologies
agree. On an intrinsic line orbit, C.28 puts every
fixed finite-word parameter track in a compact
intrinsic interval, embedded as a compact arc in
the plane. More generally, a compact intrinsic
interval $J\subset B$ lies in one of C.28's
finite-word compact sets: choose its two endpoints
there, and use the unique connecting arc in that
set. Applying a word of fixed length to every
point of $J$ remains in one larger finite-word
compact set, hence in one compact embedded
intrinsic interval. The ambient joint evaluation
map therefore gives joint intrinsic evaluation
on the path-parameter interval times $J$.
The same argument applies to the inverse word.
Consequently each path in (C.65.2) restricts to
a continuous path in the actual intrinsic
homeomorphism group of $B$, with inverse
convergence on compact intervals. All original
generators preserve the intrinsic orientation,
since their paths start at the identity; the
same is true of every word.

It follows in particular that $\rho_B(N)$ is
path connected in that intrinsic group: the
whole restricting path stays in $\rho_B(N)$
by the preceding reparametrization. This argument
uses the exact finite-word paths. It makes no
assertion about intrinsic convergence for an
arbitrary globally convergent sequence of
unbounded words.

\emph{The actual boundary image of the kernel is abelian.}
C.33 supplies a principal product chart
$O\cong A\times B_0$, where $A=C_O^0$ is the
additive line or circle. An element of
$\ker\beta$ commutes with all translations
by $A$, so its restriction has exactly the form
\[
                 (a,b)\longmapsto(a+\varphi(b),b).
                                                               \tag{C.65.5}
\]
Two such restrictions commute, including when
$A$ is a circle. Their actual global commutator
is therefore identity on $O$. Continuity makes
it identity on $\overline O$, and hence on $B$.
Thus $\rho_B(N)$ is abelian. It is normal in
$\rho_B(G)$ by the already proved normality of
$N$. No abstract source kernel has been
identified with a boundary kernel; the
commutator equality on $O$ and its extension
to $\overline O$ are the exact comparison.

\emph{A nontrivial path-connected abelian normal image forces
a metabelian action.}
Put $T=\rho_B(N)$ and $K=\rho_B(G)$, and
suppose $T\ne1$. The common fixed set of
$T$ is $K$-invariant because $T$ is normal.
Transitivity of $K$ on $B$ makes this fixed
set either all of $B$ or empty. It cannot
be all of $B$ when $T\ne1$, so it is empty.

Each $T$-orbit is path connected and nontrivial.
A nonconstant path in a line or circle contains
a nonempty open arc. Thus each orbit contains
such an arc. If an interior point of that
arc is the translate of the starting point
by an element of $T$, its inverse carries
the arc to a neighborhood of the starting
point within the same orbit. Hence every
$T$-orbit is open. The disjoint open orbits
partition the connected one-manifold $B$,
so $T$ is transitive.

An abelian transitive group acts freely:
an element fixing one point fixes every
translate of that point. Let $\overline T$
be the closure in the actual intrinsic
homeomorphism group. It is still abelian,
connected and transitive, and the same
argument makes it free. Its full centralizer
is exactly $\overline T$: a commuting map
agrees at one point with some member of the
transitive group, and that agreement
propagates to every point.

The evaluation argument of C.30, applied
in dimension one exactly as in C.48,
therefore identifies $\overline T$
homeomorphically with $B$ and makes it
locally compact. Its faithful product
action on $B\times\R^2$ is a continuous
action on a connected three-manifold.
The Hilbert--Smith input already verified
in C.33 makes it a Lie group.
It is connected of dimension one, hence
is the additive line or circle. Evaluation
is the actual conjugacy of its action
with all translations or all rotations,
respectively.

Every element of $K$ normalizes this
closure: it normalizes $T$, and conjugation
is continuous in the intrinsic
homeomorphism-group topology. On a line,
write the retained translations as $T_a$.
For a normalizing homeomorphism $g$,
\[
                 gT_ag^{-1}=T_{\theta_g(a)}
                                                               \tag{C.65.6}
\]
defines a continuous additive automorphism
of $\R$, so $\theta_g(a)=\lambda a$.
Its orientation gives $\lambda>0$.
Evaluation at zero yields
$g(a)=\lambda a+g(0)$.
Thus $K$ is a subgroup of the positive
affine group. On a circle the same
evaluation identity identifies the
induced continuous circle-group
automorphism with the degree of $g$.
That degree is $+1$, so the automorphism
is the identity and $g$ is a rotation.
In either case $K$ is metabelian.
This contradicts the hypothesis in
(C.65.3). We conclude that $T=1$,
which proves that equation.

For the claimed path comparison take
$h_s=p_sq_s^{-1}$. It is a finite-word
path starting at identity and satisfies
$\beta(h_s)=\id$ as an equality of whole
base homeomorphisms. Each $h_s$ belongs
to $N$; applying (C.65.3) gives
$\rho_B(p_s)=\rho_B(q_s)$ for every $s$.
Equality only at one varying base point
would not place this path in $\ker\beta$
and is not used here.

\emph{The full kernel and its exact discrete ambiguity.}
Let $k\in K_0$ and $g\in G$. Express $g$ by
one finite original word and replace each of its
time parameters $a_j$ by $s a_j$. This gives
an actual finite-word path $g_s$ from identity
to $g$. The path
\[
               [g_s,k]=g_s k g_s^{-1}k^{-1}
                                                               \tag{C.65.10}
\]
starts at identity, is again a single finite-word
path, and has principal-base image identity
at every time because $\beta(k)=1$.
Thus every value of (C.65.10) lies in $N$.
Equation (C.65.3) gives
$[\rho_B(g),\rho_B(k)]=1$.
Consequently $D_B$ is central in $\rho_B(G)$
and lies in its full intrinsic centralizer.

The action $\rho_B(G)$ is transitive and
nonmetabelian, hence nonabelian. C.38 identifies
that full centralizer, in its actual intrinsic
compact-open topology, as a closed discrete
group which is trivial or infinite cyclic
on a line, and finite cyclic on a circle.
Every subgroup of this closed discrete group
is closed in the intrinsic homeomorphism
group. This proves the asserted closedness,
discreteness, and cyclic description of $D_B$.
It does not require a topology on the
abstract kernel $K_0$.

Projection onto the first coordinate of
$E_B$ is a continuous surjective
homomorphism
\[
       p_B:E_B\longrightarrow\beta(G),
       \qquad \ker p_B=\Delta_B.                       \tag{C.65.11}
\]
Here both image topologies are the induced
topologies specified in the statement.
Since $D_B$ is closed and discrete in the
second factor, $\Delta_B$ is closed and
discrete in $E_B$. The preceding
centralization makes it central in $E_B$.

Write $q_B:E_B\to E_B/\Delta_B$ for the
topological group quotient. This map is
open: the inverse image of the image of
an open set $V$ is the open set
$V\Delta_B$. Choose an identity
neighborhood $U$ in $E_B$ with
$U^{-1}U\cap\Delta_B=\{1\}$, using
discreteness and continuity of multiplication
and inversion. The sets $U\delta$,
$\delta\in\Delta_B$, are pairwise disjoint,
and
\[
 q_B^{-1}(q_B(U))
       =\coprod_{\delta\in\Delta_B}U\delta.             \tag{C.65.12}
\]
The restriction of $q_B$ to each displayed
open set is a homeomorphism onto $q_B(U)$:
it is continuous and bijective, and its
openness follows from the just proved
openness of the quotient map. Translates
of these neighborhoods prove the covering
assertion.

The space $E_B$ is path connected. Indeed
the finite-word path $g_s$ used above is
continuous in both displayed image groups.
For the frontier factor this was proved
with the compact intrinsic tracks. For
the principal factor use the actual local
sections of $\pi_0$ in the product chart:
evaluation on such a section, followed
by $\pi_0$, is jointly continuous on
compact base sets, and the inverse word
gives the same fact for inverses.
This proves the required compact-open
path continuity in $\operatorname{Homeo}(B_0)$.
A deck transformation of a connected
covering is determined by its value at
one point: the agreement set of two
deck transformations is open and closed,
by evenly covered neighborhoods. Right
translation by $\Delta_B$ is transitive
on each fiber. It follows that the
entire deck group of $q_B$ is exactly
$\Delta_B$.

Finally (C.65.11) induces the continuous
bijective homomorphism
$\overline p_B:E_B/\Delta_B\to\beta(G)$.
Continuity is exactly the universal
property of the quotient topology.
The topological comparison stops at
this continuous bijection: no open-map
or inverse-continuity assertion for
$\overline p_B$ has been used or proved.

\emph{A filling inside the original domain and a prescribed
ambient disk.}
Let $\gamma$ be a based loop in $U$
which is null-homotopic in $V$.
Its compact image has an open
neighborhood in $U$. A sufficiently
fine polygonal approximation $\gamma'$
and the straight homotopies on the
small parameter intervals stay there,
fixing the basepoint. Choose finitely
many vertices generically if necessary,
so that its image is a finite planar
graph.

Take a small compact connected planar
regular neighborhood $M$ of this graph
inside $U$. For example, use small
vertex disks and thin edge rectangles
and then round their corners. Its
boundary has one outer Jordan curve
and finitely many inner Jordan curves.
Write $D_i$ for its bounded
complementary Jordan disks.
The whole closed outer Jordan disk
lies in $D$. Indeed its bounded
interior lies in the convex hull of
its boundary: a point outside that
convex hull is separated by a line
and can be joined to infinity without
meeting the boundary. The convex
hull of a compact subset of the
open Euclidean disk $D$ remains
inside $D$. In particular every
$D_i$ lies in $D$.

Fill into $M$ precisely the inner
disks $D_i$ which are entirely
contained in $V$, and call the
resulting compact planar surface $M'$.
Their boundaries already lie in $M$,
so the entire closed disk is added.
We have the actual inclusions
\[
              \gamma'\subset M\subset M'\subset V\cap D.
                                                               \tag{C.65.7}
\]
The set $M'$ is connected and meets
$U$, hence lies in that same component.
For each remaining inner disk choose
one point $z_i\in D_i\setminus V$.

The actual inclusion
\[
                  M'\hookrightarrow\R^2\setminus\{z_i\}
                                                               \tag{C.65.8}
\]
is a homotopy equivalence with an
explicit strong deformation retraction.
For each remaining closed inner Jordan
disk choose a Schoenflies disk coordinate
taking $z_i$ to zero and its boundary
to the unit circle. On the punctured
disk use
\[
        (r,\theta)\longmapsto((1-\tau)r+\tau,\theta),
                   \qquad 0<r\leq1.                    \tag{C.65.9}
\]
On the exterior of the outer Jordan
disk use its exterior Schoenflies
coordinate and the same formula for
$r\geq1$. All these homotopies fix
their boundary circles pointwise.
With the identity on $M'$, they
paste on the finite closed cover
of the punctured plane to a strong
deformation retraction onto $M'$.
Every basepoint in $M'$ stays fixed.

Since each $z_i$ lies outside $V$,
the given null-homotopy of $\gamma'$
in $V$ is a null-homotopy in the
punctured plane of (C.65.8).
Applying the displayed retraction
gives a based filling in $M'$.
Prepending the based polygonal
approximation homotopy gives a
filling of $\gamma$ in $U$.
This proves (C.65.4).

For original ambient-null evaluation loops
converging to a $G$-fixed frontier
point, compactness of the original
time interval makes their entire
tracks shrink to that point.
Choose open Euclidean disks
containing those tracks with
radii tending to zero. Each track
lies in one connected component
of its original orbit intersected
with that disk, and C.32 supplies
its null-homotopy in the original
orbit. The just proved injection
gives a filling in that same
component, with the stated
vanishing ambient diameter.

The first part compares specified
finite-word kernel paths and records
the central discrete ambiguity of the
full kernel through an actual covering.
It does not construct a new base lift
or identify the final continuous
bijection with a homeomorphism.
It also does not factor arbitrary
source germs through a frontier action.
The second part chooses each
filling separately; it does not
give a continuous choice of
fillings, simultaneous configuration
or edge isotopy extension, or
fragmentation of an arbitrary
original loop. These remaining
requirements of C.12 are retained.
\end{proof}
